\documentclass[12pt,a4paper]{report}

\usepackage[utf8]{inputenc}
\usepackage[T1]{fontenc}
\usepackage{lmodern}
\usepackage{geometry}
\usepackage{setspace}
\usepackage{amsmath,amssymb,amsthm}
\usepackage{graphicx}
\usepackage{tikz}
\usetikzlibrary{arrows.meta,positioning,calc}
\usepackage{subcaption}
\usepackage{pgfplots}
\usepackage{pgfplotstable}
\usepackage{booktabs}
\usepackage{longtable}
\usepackage{algorithm}
\usepackage{algpseudocode}
\usepackage{hyperref}
\usepackage{caption}
\usepackage{enumitem}
\usepackage[strings]{underscore}

\pgfplotsset{compat=1.18}
\pgfplotstableset{col sep=comma}

\geometry{margin=1in}
\onehalfspacing

\newtheorem{theorem}{Theorem}[chapter]
\newtheorem{lemma}[theorem]{Lemma}
\newtheorem{proposition}[theorem]{Proposition}
\newtheorem{corollary}[theorem]{Corollary}
\theoremstyle{definition}
\newtheorem{definition}[theorem]{Definition}
\newtheorem{example}[theorem]{Example}
\theoremstyle{remark}
\newtheorem{remark}[theorem]{Remark}

\newcommand{\Graph}{G}
\newcommand{\Face}{F}
\newcommand{\ChordSet}{S}
\newcommand{\Root}{r}
\newcommand{\GS}{\mathrm{GS}}
\newcommand{\VGS}{\mathrm{VGS}}

% Point to the results CSV relative to this file's location.
% Adjust the path if your directory layout differs.
\pgfplotstableread[col sep=comma]{results.csv}\resultsdata

\begin{document}

% -----------------------------------------------------------------------
%  Title page
% -----------------------------------------------------------------------
\begin{titlepage}
    \centering
    {\Large Bangladesh University of Engineering and Technology (BUET)\par}
    \vspace{1.0cm}
    {\LARGE Department of Computer Science and Engineering\par}
    \vspace{1.5cm}
    {\LARGE \textbf{Generating All Triangulations of Biconnected Plane Graphs\\[4pt]
                     in Amortized Constant Time per Triangulation}\par}
    \vspace{1.5cm}
    {\large Undergraduate Thesis\par}
    \vspace{1.5cm}
    {\large Tawkir Aziz Rahman\par}
    \vspace{0.5cm}
    {\large Dhaka-1000, Bangladesh\par}
    \vfill
    {\large \today\par}
\end{titlepage}

% -----------------------------------------------------------------------
%  Front matter
% -----------------------------------------------------------------------
\pagenumbering{roman}

\chapter*{Abstract}
\addcontentsline{toc}{chapter}{Abstract}

We present an algorithm for generating all triangulations of a biconnected plane graph
$\Graph$ with $n$ vertices in $O(1)$ amortized time per triangulation and $O(n)$ space.
The method builds a two-layer structure: a local genealogical tree of triangulations for
each face and a global depth-first search (DFS) that combines those face triangulations
without creating multi-edges. The key idea is a safe root vertex selection that prevents
conflicts across faces, together with a pruning strategy and a valid generating set that
keep each flip constant-time. The thesis revisits the cycle triangulation algorithm of
Parvez, Rahman, and Nakano~\cite{parvez2011}, identifies why it fails on biconnected
graphs, and then shows how to extend the approach while preserving the amortized $O(1)$
bound. We also provide empirical validation using the measurements recorded in
\texttt{results.csv}.

\tableofcontents
\listoffigures
\listofalgorithms
\clearpage

\pagenumbering{arabic}

% -----------------------------------------------------------------------
%  Chapter 1 – Introduction
% -----------------------------------------------------------------------
\chapter{Introduction}
\label{chap:introduction}

To achieve our objective of generating all triangulations for biconnected plane graphs,
we first study the closest previous approaches. The most relevant work is the 2011 JGAA
paper by Parvez, Rahman, and Nakano~\cite{parvez2011}. In that work, the authors present
an algorithm that generates all triangulations of a cycle in amortized constant time per
triangulation using $O(n)$ space, where $n$ is the number of vertices in the cycle. Their
main idea is to organize triangulations into a genealogical tree whose nodes are canonical
triangulations, and to define a unique parent-child relationship by a local flip rule.

That cycle algorithm is one of the main parts of our method. However, our target is
broader: we want to generate all triangulations of biconnected plane graphs. The
triconnected case in~\cite{parvez2011} can treat faces independently because multi-edges
cannot occur. In biconnected graphs, separation pairs appear and allow the same chord to
arise in two different faces, creating the multi-edge problem. Therefore, the main
technical challenge is to preserve the constant-time flip structure while preventing
multi-edges across faces.

In this thesis, we keep the author's algorithmic intuition and expand it into a full
thesis structure with formal definitions, proofs, and experiments. The flow is as
follows. Chapter~\ref{chap:preliminaries} provides preliminaries.
Chapter~\ref{chap:cycle} revisits the cycle triangulation algorithm.
Chapter~\ref{chap:gsdistinction} explains the distinction between the original generating
chord notion and our generating set. Chapters~\ref{chap:triconnected}
and~\ref{chap:faces} describe why the triconnected algorithm fails for biconnected graphs
and introduce the face-by-face DFS approach. Chapters~\ref{chap:pruning}
and~\ref{chap:saferoot} develop the pruning strategy and safe root theory.
Chapter~\ref{chap:vgs} introduces the valid generating set and constant-time updates.
Chapter~\ref{chap:algorithm} presents the full algorithm.
Chapter~\ref{chap:complexity} analyzes time and space complexity,
Chapter~\ref{chap:experiments} reports experiments from \texttt{results.csv},
Chapter~\ref{chap:limitations} discusses one-connected graphs,
and Chapter~\ref{chap:conclusion} concludes.

% -----------------------------------------------------------------------
%  Chapter 2 – Preliminaries
% -----------------------------------------------------------------------
\chapter{Preliminaries}
\label{chap:preliminaries}

We define the terms used throughout the thesis. Let $\Graph=(V,E)$ be a connected
simple graph with vertex set $V$ and edge set $E$. An edge connecting vertices $v_i$
and $v_j$ is denoted by $(v_i,v_j)$ where $v_i \neq v_j$. The connectivity of a graph
is the minimum number of vertices whose removal disconnects the graph. A graph is
\textit{biconnected} if it remains connected after removing any single vertex, and
\textit{triconnected} if it remains connected after removing any two
vertices~\cite{diestel2017,hopcroft1974}.

A graph is \textit{planar} if it can be embedded in the plane so that no edges cross
except at their endpoints. A \textit{plane graph} is a planar graph with a fixed
embedding. This embedding divides the plane into regions called faces. The unbounded
region is the outer face; all other regions are inner faces. A plane graph is a
\textit{plane triangulation} if each face is a triangle. Note that edges are not required
to be straight line segments.

A \textit{cycle} $C$ in a plane graph is a sequence of distinct vertices
$v_1,v_2,\ldots,v_k$ such that $(v_i,v_{i+1})\in E$ for $1 \le i < k$ and
$(v_k,v_1)\in E$. A face boundary is such a cycle. A \textit{chord} of a cycle is an
edge between two non-consecutive vertices that lies inside the face. A set of
non-intersecting chords that decomposes the cycle into triangles is a \textit{triangulation}
of the cycle.

\begin{figure}[h]
    \centering
    \begin{subfigure}{0.45\textwidth}
        \centering
        \begin{tikzpicture}[scale=1, every node/.style={circle, draw, fill=white, inner sep=2pt}]
            \foreach \i/\angle in {1/90, 2/30, 3/-30, 4/-90, 5/-150, 6/150} {
                \node (a\i) at (\angle:1) {\i};
            }
            \draw (a1) -- (a2) -- (a3) -- (a4) -- (a5) -- (a6) -- (a1);
            \draw (a2) -- (a4) -- (a6) -- (a2);
        \end{tikzpicture}
        \caption{Triangulation A}
    \end{subfigure}
    \hfill
    \begin{subfigure}{0.45\textwidth}
        \centering
        \begin{tikzpicture}[scale=1, every node/.style={circle, draw, fill=white, inner sep=2pt}]
            \foreach \i/\angle in {1/90, 2/30, 3/-30, 4/-90, 5/-150, 6/150} {
                \node (b\i) at (\angle:1) {\i};
            }
            \draw (b1) -- (b2) -- (b3) -- (b4) -- (b5) -- (b6) -- (b1);
            \draw (b1) -- (b3) -- (b5) -- (b1);
        \end{tikzpicture}
        \caption{Triangulation B}
    \end{subfigure}
    \caption{Two distinct triangulations of a cycle}
    \label{fig:distinct_triangulations}
\end{figure}

In a triangulation $T$ of a cycle, the set of chords is maximal: any chord not in $T$
would intersect a chord in $T$. We say a vertex $v_j$ is \textit{visible} from $v_i$ if
the chord $(v_i,v_j)$ is present in the triangulation.

Throughout this thesis, we work with labeled graphs, so each vertex has a unique label.
This assumption matches the algorithmic setting of~\cite{parvez2011} and the experimental
data in \texttt{results.csv}.

% -----------------------------------------------------------------------
%  Chapter 3 – Triangulating a Cycle in Amortized Constant Time
% -----------------------------------------------------------------------
\chapter{Triangulating a Cycle in Amortized Constant Time}
\label{chap:cycle}

We first revisit the cycle triangulation algorithm from~\cite{parvez2011} because it is
the backbone of our approach. Let $C$ be a cycle with vertices $v_1,v_2,\ldots,v_n$ in
counterclockwise order. We build a genealogical tree of triangulations in which each node
is a triangulation and each edge corresponds to a single chord flip. The parent-child
relationship is defined using blocking chords and a leftmost rule, which yields a unique
path from the root triangulation to every other triangulation.

\section{Root vertex and visibility}

We pick a root vertex, say $v_1$, and classify chords according to visibility from $v_1$.
If every non-neighboring vertex is connected to $v_1$, then all vertices are visible and
the triangulation is a root triangulation.

\begin{definition}[Blocking chord]
A chord $(v_i,v_j)$ is a \textit{blocking chord} with respect to root $v_1$ if it does
not have $v_1$ as an endpoint and is visible from $v_1$, meaning it forms a triangle with
$v_1$ and blocks visibility to some vertices behind it.
\end{definition}

\begin{definition}[Generating chord]
A chord incident to the root $v_1$ is called a \textit{generating chord}. In the cycle
algorithm of~\cite{parvez2011}, a generating chord is flippable if its flip becomes the
leftmost blocking chord of the child triangulation.
\end{definition}

\section{Chord flipping}

A flip is defined on a quadrilateral. Suppose a triangulation contains a chord $(b,d)$
inside a quadrilateral $a,b,c,d$ in order around the boundary. The flip removes $(b,d)$
and adds the other diagonal $(a,c)$. This is a constant-time operation and changes only
one chord.

\begin{figure}[h]
    \centering
    \begin{tikzpicture}[scale=1.1, every node/.style={circle, draw, fill=white, inner sep=2pt}]
        \node (a) at (0,0) {$a$};
        \node (b) at (2,0) {$b$};
        \node (c) at (2,2) {$c$};
        \node (d) at (0,2) {$d$};
        \draw (a) -- (b) -- (c) -- (d) -- (a);
        \draw[thick, red] (b) -- (d);
        \draw[thick, blue, dashed] (a) -- (c);
    \end{tikzpicture}
    \caption{Flipping chord $(b,d)$ to $(a,c)$ in a quadrilateral}
    \label{fig:flip}
\end{figure}

\section{Leftmost blocking chord}

To ensure a unique parent,~\cite{parvez2011} defines the \textit{leftmost blocking
chord}. Intuitively, among all blocking chords, the one that blocks the leftmost blocked
vertex is selected. The parent of a triangulation is obtained by flipping this leftmost
blocking chord.

\begin{lemma}[Uniqueness of leftmost blocking chord]
If a triangulation has at least one blocking chord, then it has a unique leftmost blocking
chord.
\end{lemma}

\begin{proof}
Suppose there are two distinct blocking chords that both block the leftmost blocked
vertex. The root vertex must see each blocking chord to classify it as blocking, so each
blocking chord forms a triangle with the root. If two blocking chords both block the same
leftmost vertex, then one of them must lie closer to the root. That chord blocks
visibility to the other, contradicting that both are visible from the root. Therefore,
the leftmost blocking chord is unique.
\end{proof}

\section{Parent-child relationship}

A triangulation is the parent of another if flipping the child's leftmost blocking chord
produces the parent. By the uniqueness above, each triangulation has exactly one parent
(except the root). Thus the genealogical tree has a unique path from the root to any
node.

When we keep flipping leftmost blocking chords, we eventually reach a triangulation with
no blocking chord, meaning all chords are incident to the root. This is the root
triangulation.

\section{Generating children}

To generate children of a triangulation, we flip generating chords. The key observation
from~\cite{parvez2011} is that only certain generating chords are eligible: if
$(v_r,v_{r'})$ is the leftmost blocking chord of a triangulation, then flipping a
generating chord $(v_1,v_j)$ can create a child whose leftmost blocking chord is the new
chord only when $j \ge r$.

\begin{lemma}[Flippable generating chords]
Let $(v_r,v_{r'})$ with $r < r'$ be the leftmost blocking chord of a triangulation $T$.
Then flipping $(v_1,v_j)$ results in a child whose new chord is the leftmost blocking
chord if and only if $j \ge r$.
\end{lemma}

\begin{proof}
If $j < r$, then the flip stays inside the region bounded by $(v_1,v_{r'})$, so the
leftmost blocking chord $(v_r,v_{r'})$ remains visible and leftmost, and the new chord
cannot become the leftmost blocking chord. If $j \ge r$, the flipped chord lies to the
left of $(v_r,v_{r'})$ in the visibility order, so the new chord is visible from the root
and blocks a vertex no farther right than $v_{r'}$. Therefore, it becomes the leftmost
blocking chord of the child.
\end{proof}

\section{Data structures and complexity}

The cycle algorithm stores the triangulation using three sets:
\begin{itemize}[leftmargin=2em]
    \item $\GS$: generating chords incident to the root.
    \item $O$: opposite pairs used to perform constant-time flips.
    \item $L$: all remaining chords of the triangulation.
\end{itemize}
Since the cycle has $n$ vertices, each set is $O(n)$ in size. Each flip changes only a
constant number of chords, so moving to a child triangulation takes $O(1)$ time. The
depth of any branch is at most $n-2$, so recursion uses $O(n)$ stack space. Therefore,
the algorithm runs in amortized $O(1)$ time per triangulation with $O(n)$ space.

\section{Cycle triangulation algorithm}

Algorithm~\ref{alg:cycle} summarizes the cycle enumeration procedure used in the
triconnected case and adapted later in this thesis.

\begin{algorithm}[H]
\caption{EnumerateCycleTriangulations}
\label{alg:cycle}
\begin{algorithmic}[1]
\Procedure{EnumerateCycleTriangulations}{$C$}
    \State Choose root vertex $v_1$
    \State Build root triangulation by adding all chords $(v_1,v_j)$ for non-neighboring $v_j$
    \State Initialize $\GS$, $O$, and $L$
    \State \Call{DFSLocal}{$T_{root}$}
\EndProcedure
\Statex
\Procedure{DFSLocal}{$T$}
    \State Output $T$
    \If{$T$ has a leftmost blocking chord $(v_r,v_{r'})$}
        \For{each generating chord $(v_1,v_j)$ with $j \ge r$}
            \State Flip $(v_1,v_j)$ to obtain $T'$
            \State \Call{DFSLocal}{$T'$}
        \EndFor
    \EndIf
\EndProcedure
\end{algorithmic}
\end{algorithm}

% -----------------------------------------------------------------------
%  Chapter 4 – Distinction in Generating Set Definition
% -----------------------------------------------------------------------
\chapter{Distinction in Generating Set Definition}
\label{chap:gsdistinction}

The cycle algorithm in~\cite{parvez2011} calls only those chords flippable if their flip
becomes the leftmost blocking chord. In our thesis, we distinguish two sets:
\begin{itemize}[leftmargin=2em]
    \item $\GS$: all chords incident to the root (every chord with one endpoint at the root).
    \item $\VGS$: the valid subset of $\GS$ whose flips do not create multi-edges with
          previously processed faces.
\end{itemize}
This distinction is necessary because in biconnected graphs a generating chord may lead
to a multi-edge in another face. We therefore maintain a separate valid generating set
that is updated after each flip.

% -----------------------------------------------------------------------
%  Chapter 5 – Why the Triconnected Algorithm Fails for Biconnected Graphs
% -----------------------------------------------------------------------
\chapter{Why the Triconnected Algorithm Fails for Biconnected Graphs}
\label{chap:triconnected}

Parvez, Rahman, and Nakano~\cite{parvez2011} extend the cycle algorithm to triconnected
plane graphs by triangulating all faces independently. The key reason this works is that
triconnected graphs have no separation pairs, so a chord added in one face cannot also
appear in another face. Consequently, no multi-edge can occur.

In biconnected graphs, separation pairs exist. A separation pair $(a,b)$ splits the graph
into two components if removed. The two components can each place a chord between $a$ and
$b$ on different faces. If we triangulate those faces independently, the chord $(a,b)$ may
appear twice, creating a multi-edge. This invalidates the assumption of independent face
triangulation.

\begin{figure}[h]
    \centering
    \begin{tikzpicture}[scale=0.7, line cap=round, line join=round]
        \tikzset{v/.style={circle,draw,fill=white,inner sep=1.8pt,minimum size=14pt}}
        \coordinate (L1) at (0,2);
        \coordinate (L2) at (0,6);
        \coordinate (L3) at (2,6);
        \coordinate (L4) at (4,6);
        \coordinate (L5) at (4,4);
        \coordinate (L6) at (4,2);
        \coordinate (L7) at (2,2);
        \coordinate (L8) at (2,4);
        \draw (L1)--(L2)--(L3)--(L4)--(L5)--(L6)--(L7)--(L8);
        \draw (L1)--(L7);
        \draw (L3)--(L8);
        \node[v] at (L1) {1};
        \node[v] at (L2) {2};
        \node[v] at (L3) {3};
        \node[v] at (L4) {4};
        \node[v] at (L5) {5};
        \node[v] at (L6) {6};
        \node[v] at (L7) {7};
        \node[v] at (L8) {8};
    \end{tikzpicture}
    \caption{A biconnected graph where $(3,7)$ is a separation pair. The chord $(3,7)$ can appear in two faces.}
    \label{fig:separation_pair}
\end{figure}

This multi-edge issue is the core reason we cannot directly apply the triconnected
algorithm to biconnected graphs. We must process faces in an order that respects chords
already chosen, and we must prune invalid branches early.

% -----------------------------------------------------------------------
%  Chapter 6 – Processing Faces One by One
% -----------------------------------------------------------------------
\chapter{Processing Faces One by One}
\label{chap:faces}

To avoid multi-edges, we process faces one by one using a DFS across faces. For each
face $F_i$, we generate its local genealogical tree and traverse it. For every
triangulation of $F_i$, we recursively process the next face, carrying a global chord
set $\ChordSet$ that tracks all chords chosen so far.

Consider faces $F_1, F_2, \ldots, F_k$. For each triangulation $T_{1,j}$ of $F_1$, we
enumerate all triangulations $T_{2,*}$ of $F_2$ that do not conflict with $T_{1,j}$.
For each compatible pair, we move to $F_3$, and so on, until we triangulate $F_k$. Each
root-to-leaf path corresponds to a complete triangulation of $\Graph$.

\begin{figure}[h]
    \centering
    \begin{tikzpicture}[
        level distance=1.2cm,
        sibling distance=2.2cm,
        every node/.style={rectangle, rounded corners, draw, minimum width=1.8cm,
                           minimum height=0.7cm, align=center, font=\small},
        edge from parent/.style={draw, -latex}
    ]
    \node {Face $F_1$}
        child { node {Face $F_2$}
            child { node {Face $F_3$} }
            child { node {Face $F_3$} }
        }
        child { node {Face $F_2$}
            child { node {Face $F_3$} }
        };
    \end{tikzpicture}
    \caption{DFS traversal across faces. Each branch corresponds to one compatible choice per face.}
    \label{fig:face_dfs}
\end{figure}

This layered DFS preserves the author's intended flow: we keep the local genealogical
trees from the cycle algorithm but combine them in a global search that respects already
chosen chords.

% -----------------------------------------------------------------------
%  Chapter 7 – Pruning Strategy and Conflict Tracking
% -----------------------------------------------------------------------
\chapter{Pruning Strategy and Conflict Tracking}
\label{chap:pruning}

A naive approach would generate all local triangulations and then filter invalid ones.
This is too expensive because invalid triangulations could dominate the search and destroy
the amortized constant-time bound. Instead, we prune branches as soon as a conflicting
chord appears.

\section{Persistence of non-root chords}

In the cycle algorithm, we only flip generating chords (those incident to the root).
Therefore, once a chord is introduced that is not incident to the root, it is never
flipped again in the descendant triangulations.

\begin{lemma}[Chord persistence]
Let $T$ be a triangulation of a face with root vertex $v_1$. If a chord $(a,b)$ does not
have $v_1$ as an endpoint, then $(a,b)$ appears in every descendant triangulation of $T$
in the local genealogical tree.
\end{lemma}

\begin{proof}
Descendants are obtained only by flipping generating chords incident to $v_1$. A chord
$(a,b)$ with neither endpoint $v_1$ is never selected for flipping, and thus cannot be
removed. Therefore it persists in all descendants.
\end{proof}

This property justifies pruning: if a newly flipped chord creates a multi-edge with any
chord already in $\ChordSet$, then every descendant triangulation will also contain this
conflicting chord and therefore will be invalid. We can safely prune that branch without
losing any valid triangulation.

\section{Global chord set}

We maintain a global chord set $\ChordSet$ (implemented as a hash map) that stores all
chords from previously processed faces and the permanent edges of the graph. Each time we
move to a new triangulation of face $F_i$, we insert its chords into $\ChordSet$. When we
backtrack, we remove them. This supports constant-time conflict checks for each flip.

% -----------------------------------------------------------------------
%  Chapter 8 – Safe Root Vertices
% -----------------------------------------------------------------------
\chapter{Safe Root Vertices}
\label{chap:saferoot}

The pruning strategy assumes that the root triangulation of each face is valid. Therefore,
we must guarantee that every face has a \textit{safe root} vertex such that the fan
triangulation from that root does not conflict with $\ChordSet$.

\section{Existence of safe roots}

\begin{theorem}[Existence of safe roots]
\label{thm:safe_root}
For any face $F_i$ with $n \ge 3$ vertices in a biconnected plane graph, and any global
chord set $\ChordSet$ from previously processed faces, there exist at least two safe root
vertices in $F_i$.
\end{theorem}

\begin{proof}
If no chord in $\ChordSet$ has both endpoints on the boundary of $F_i$, then every vertex
is safe. Otherwise, let $(v_i,v_j)$ be a chord in $\ChordSet$ with endpoints on the
boundary. By planarity, this chord separates the boundary into two disjoint subpaths.
Any external chord must have both endpoints inside the same subpath; chords cannot cross
$(v_i,v_j)$. Therefore, the boundary is partitioned into two structures, each of which is
a path with an external chord between its endpoints. We call each such subpath a $T_n$
structure.

For $T_1$, the middle vertex has no non-neighboring vertex in the structure, so it is
safe. For $T_2$, if one interior vertex has an external chord, the other interior vertex
becomes a $T_1$ case and is safe. By induction on $n$, each $T_n$ structure has at least
one safe vertex. Since there are at least two such structures, there are at least two safe
roots in $F_i$.
\end{proof}

\section{Minimum number of valid triangulations}

The safe root theorem implies at least two distinct fan triangulations for each face.
Those two root triangulations differ in at least $n-4$ chords, so the genealogical tree
contains at least $n-3$ distinct triangulations reachable by flips. Therefore, each face
has at least $\Omega(n)$ valid triangulations, which is enough to amortize the $O(n)$
time for building the root triangulation and valid generating set.

\section{Finding a safe root in $O(n)$ time}

We can find a safe root using a two-pointer scan along the boundary. Let the vertices be
$v_1,\ldots,v_n$ in order. We initialize $i=1$ and $j=n-1$ (neighbors are excluded). If
$(v_i,v_j)\in\ChordSet$, then $v_i$ cannot be a safe root and we advance $i$. Otherwise,
we decrease $j$ because $v_i$ still may be safe and we need to test shorter chords. This
scan ends when $i+1=j$, at which point $v_i$ is safe.

\begin{algorithm}[H]
\caption{FindSafeRoot}
\label{alg:saferoot}
\begin{algorithmic}[1]
\Procedure{FindSafeRoot}{$F_i,\ChordSet$}
    \State Let $v_1,\ldots,v_n$ be the boundary of $F_i$
    \State $i \gets 1$, $j \gets n-1$
    \While{$j > i+1$}
        \If{$(v_i,v_j) \in \ChordSet$}
            \State $i \gets i+1$
        \Else
            \State $j \gets j-1$
        \EndIf
    \EndWhile
    \State \Return $v_i$
\EndProcedure
\end{algorithmic}
\end{algorithm}

The loop advances at least one pointer per step, so the time is $O(n)$. This matches the
$O(n)$ cost of building the root fan triangulation and preserves the amortized
constant-time bound.

% -----------------------------------------------------------------------
%  Chapter 9 – Valid Generating Set and Constant-Time Updates
% -----------------------------------------------------------------------
\chapter{Valid Generating Set and Constant-Time Updates}
\label{chap:vgs}

We now introduce the valid generating set ($\VGS$), which contains exactly those
generating chords whose flips do not create multi-edges with $\ChordSet$. During the
initial root construction we compute $\VGS$ by checking each generating chord's opposite
pair against $\ChordSet$.

After each flip, only the opposite pairs of the four boundary edges of the quadrilateral
change. Therefore, we only need to update the validity of those four chords. There are
four cases:
\begin{enumerate}[leftmargin=2em]
    \item The chord was valid and becomes invalid after the flip (remove from $\VGS$).
    \item The chord was invalid and becomes valid after the flip (insert into $\VGS$).
    \item The chord remains valid.
    \item The chord remains invalid.
\end{enumerate}
If $\VGS$ is stored as a linked list (or a hash set with pointers), each insert or delete
is $O(1)$. Therefore, each flip still runs in $O(1)$ time while maintaining a correct
set of valid generating chords.

% -----------------------------------------------------------------------
%  Chapter 10 – Full Algorithm
% -----------------------------------------------------------------------
\chapter{Full Algorithm}
\label{chap:algorithm}

We now present the full algorithm in pseudocode. The main procedure processes faces in
order, using the safe root algorithm for each face and enumerating valid triangulations
with the local genealogical tree and $\VGS$ updates.

\begin{algorithm}[H]
\caption{FindAllTriangulations}
\label{alg:full}
\begin{algorithmic}[1]
\Procedure{FindAllTriangulations}{$\Graph$}
    \State Label faces $F_1, F_2, \ldots, F_k$
    \State Initialize global chord set $\ChordSet$ with all permanent edges
    \State Initialize empty partial triangulation $T$
    \State \Call{ProcessFace}{$1, T, \ChordSet$}
\EndProcedure
\end{algorithmic}
\end{algorithm}

\begin{algorithm}[H]
\caption{ProcessFace}
\label{alg:processface}
\begin{algorithmic}[1]
\Procedure{ProcessFace}{$i, T, \ChordSet$}
    \If{$i > k$}
        \State Output $T$
        \State \Return
    \EndIf
    \State $\Root \gets$ \Call{FindSafeRoot}{$F_i, \ChordSet$}
    \State Build root triangulation $T_i$ of $F_i$ using $\Root$
    \State Initialize $\GS$, $O$, $L$, and $\VGS$ for $T_i$
    \State \Call{EnumerateFace}{$F_i, T_i, T, \ChordSet$}
\EndProcedure
\Statex
\Procedure{EnumerateFace}{$F_i, T_i, T, \ChordSet$}
    \State Add chords of $T_i$ to $\ChordSet$
    \State \Call{ProcessFace}{$i+1, T \cup T_i, \ChordSet$}
    \State Remove chords of $T_i$ from $\ChordSet$
    \ForAll{chords $e \in \VGS$}
        \State Flip $e$ to obtain $T_i'$
        \State Update opposite pairs and $\VGS$ for the four boundary edges
        \If{new chord not in $\ChordSet$}
            \State \Call{EnumerateFace}{$F_i, T_i', T, \ChordSet$}
        \EndIf
    \EndFor
\EndProcedure
\end{algorithmic}
\end{algorithm}

This pseudocode follows the author's flow: each face has its own genealogical tree, the
DFS across faces is the outer layer, and pruning is enforced by the global chord set.

% -----------------------------------------------------------------------
%  Chapter 11 – Complexity Analysis
% -----------------------------------------------------------------------
\chapter{Complexity Analysis}
\label{chap:complexity}

\section{Time complexity}

For each face, building the root triangulation and the initial valid generating set costs
$O(n)$. By Theorem~\ref{thm:safe_root} and the argument in
Chapter~\ref{chap:saferoot}, each face has at least $\Omega(n)$ valid triangulations, so
this $O(n)$ cost is amortized to $O(1)$ per triangulation within the face.

Each flip updates only a constant number of chords and opposite pairs, so the transition
from one triangulation to the next takes $O(1)$ time. The outer DFS across faces does
not introduce additional asymptotic overhead because each complete triangulation
corresponds to exactly one leaf in the recursion tree and is generated by a constant
amount of work per flip. Therefore, the overall algorithm runs in $O(1)$ amortized time
per triangulation.

\section{Space complexity}

We store the current face's $\GS$, $\VGS$, $O$, and $L$ sets, each of which is $O(n)$
for a face with $n$ vertices. The global chord set $\ChordSet$ stores all permanent edges
and all chords along the current recursion path, which is $O(n)$ because the number of
edges in a plane graph is linear in $n$. The recursion depth across faces is also $O(n)$
since the number of faces is linear. Therefore, the total space complexity is $O(n)$.

% -----------------------------------------------------------------------
%  Chapter 12 – Experimental Validation
% -----------------------------------------------------------------------
\chapter{Experimental Validation}
\label{chap:experiments}

This section reports empirical validation using the latest measurements in
\texttt{../time\_complexity/results.csv}. Each row records: the instance name, number of
vertices, number of triangulations, mean time, median time, minimum and maximum time,
standard deviation, per-triangulation time in nanoseconds, peak memory, and memory per
vertex. We include a sample table in the main text and provide the full dataset in
Appendix~\ref{app:full_results}.

\section{Sample summary table}

% \begin{table}[h]
% \centering
% \caption{Sample of experimental results (first 10 rows)}
% \label{tab:sample}
% \pgfplotstabletypeset[
%     columns={filename,vertices,triangulations,meanTime,perTriangNs,peakMemory},
%     columns/filename/.style={string type,column name=Instance},
%     columns/vertices/.style={column name=$|V|$},
%     columns/triangulations/.style={column name=Triangulations},
%     columns/meanTime/.style={column name=Mean Time (s), fixed, precision=6},
%     columns/perTriangNs/.style={column name=Per Triangulation (ns)},
%     columns/peakMemory/.style={column name=Peak Memory},
%     every head row/.style={before row=\toprule, after row=\midrule},
%     every last row/.style={after row=\bottomrule},
%     rows={0,...,9}
% ]{\resultsdata}
% \end{table}

\section{Per-triangulation time plot}

% \begin{figure}[h]
% \centering
% \begin{tikzpicture}
% \begin{axis}[
%     width=0.9\textwidth,
%     height=0.45\textwidth,
%     xlabel=Vertices,
%     ylabel=Per-triangulation time (ns),
%     grid=both,
%     grid style={dashed,gray!40},
%     tick label style={font=\small}
% ]
% \addplot[only marks, mark=*, mark size=1.2pt]
%     table[x=vertices,y=perTriangNs]{../time_complexity/results.csv};
% \end{axis}
% \end{tikzpicture}
% \caption{Per-triangulation time versus number of vertices across the dataset.}
% \label{fig:per_triang_time}
% \end{figure}

The measurements confirm that per-triangulation time remains small across a wide range of
instances. The full dataset is included in the appendix for reproducibility.

% -----------------------------------------------------------------------
%  Chapter 13 – Why the Algorithm Does Not Extend to 1-Connected Graphs
% -----------------------------------------------------------------------
\chapter{Why the Algorithm Does Not Extend to 1-Connected Graphs}
\label{chap:limitations}

In a 1-connected plane graph, a separation vertex can appear multiple times on the
boundary of a single face. This can create faces in which the same vertex repeats in the
boundary cycle. In such a case, no safe root may exist. For example, in the path graph
$1$-$2$-$3$-$4$, the unique face boundary is $1$-$2$-$3$-$4$-$3$-$2$-$1$. Any root
choice creates a repeated chord or a self-loop, so no safe root triangulation exists.
Hence the algorithm cannot be directly extended to all 1-connected graphs.

However, some 1-connected graphs do admit safe roots. For example, in the face
$1$-$2$-$3$-$4$-$5$-$2$-$1$, vertices $1$, $3$, and $5$ can be safe roots, and the
algorithm still works. Therefore, the limitation is not absolute; it depends on the
existence of safe roots in all faces.

% -----------------------------------------------------------------------
%  Chapter 14 – Limitations
% -----------------------------------------------------------------------
\chapter{Limitations}

The current algorithm is limited to biconnected plane graphs because the existence of
safe roots is not guaranteed in general 1-connected graphs with repeated vertices on face
boundaries. The pruning strategy also relies on the persistence of non-root chords, which
fails if a face boundary is not a simple cycle.

% -----------------------------------------------------------------------
%  Chapter 15 – Future Work
% -----------------------------------------------------------------------
\chapter{Future Work}
\label{chap:futurework}

Future work includes characterizing which 1-connected graphs admit safe roots in all
faces and extending the pruning strategy to tolerate a controlled number of invalid
generating chords. Another direction is to explore 1-planar graphs, where at most one
chord can cross another. Extending the genealogical tree approach to that setting could
lead to new enumeration algorithms.

% -----------------------------------------------------------------------
%  Chapter 16 – Conclusion
% -----------------------------------------------------------------------
\chapter{Conclusion}
\label{chap:conclusion}

We presented an algorithm that enumerates all triangulations of a biconnected plane graph
in $O(1)$ amortized time per triangulation and $O(n)$ space. The algorithm preserves the
canonical genealogical structure of the Parvez-Rahman-Nakano cycle algorithm while
resolving the multi-edge issue introduced by separation pairs. The key ideas are safe root
selection, a valid generating set, and a pruning strategy based on the persistence of
non-root chords. Empirical results from the latest dataset confirm the practical
efficiency of the approach.

% -----------------------------------------------------------------------
%  Appendix
% -----------------------------------------------------------------------
\appendix
\chapter{Full Experimental Results}
\label{app:full_results}

\small
% \pgfplotstabletypeset[
%     columns={filename,vertices,triangulations,meanTime,medianTime,minTime,maxTime,
%              stddevTime,perTriangNs,peakMemory,memoryPerVertex},
%     columns/filename/.style={string type,column name=Instance},
%     columns/vertices/.style={column name=$|V|$},
%     columns/triangulations/.style={column name=Triangulations},
%     columns/meanTime/.style={column name=Mean},
%     columns/medianTime/.style={column name=Median},
%     columns/minTime/.style={column name=Min},
%     columns/maxTime/.style={column name=Max},
%     columns/stddevTime/.style={column name=StdDev},
%     columns/perTriangNs/.style={column name=PerTriang (ns)},
%     columns/peakMemory/.style={column name=PeakMem},
%     columns/memoryPerVertex/.style={column name=Mem/Vertex},
%     begin table=\begin{longtable}{lrrrrrrrrrr},
%     end table=\end{longtable},
%     every head row/.style={before row=\toprule, after row=\midrule},
%     every last row/.style={after row=\bottomrule}
% ]{\resultsdata}

% -----------------------------------------------------------------------
%  Bibliography
% -----------------------------------------------------------------------
\bibliographystyle{plain}
\bibliography{references}

\end{document}